\section{Metodología de Investigación}

La investigación se estructuró en cinco frentes ejecutados en paralelo con consulta web el 23 de julio de 2026:

\begin{enumerate}
	\item Condiciones climáticas de diseño del sitio (IDEAM/ASHRAE/proxies aeroportuarios y verificación de densidad por gas ideal).
	\item Ventiladores para ambiente corrosivo (Greenheck, Hartzell, New York Blower, Plastec, Sodeca, Soler \& Palau, Systemair).
	\item Filtros MERV 13-14 y equivalencias ISO 16890/EN 779 (Camfil, AAF, Koch, Freudenberg y canal colombiano).
	\item Rejillas de exfiltración, dampers barométricos de alivio y mallas anti-insectos (Greenheck, Ruskin, Nailor, TROX, Halton, Titus, fabricantes nacionales).
	\item Instrumentación de presión diferencial y accesorios de montaje (Dwyer, Setra, Siemens, Honeywell).
\end{enumerate}

Como criterios de aceptación de fuentes se privilegiaron fichas y catálogos de fabricante con datos de desempeño ensayados (AMCA 210/500-D, ASHRAE 52.2, ASHRAE 70); toda cifra comercial se cita con URL y fecha de consulta. Las magnitudes sin fuente de catálogo concreta se marcan como dato típico y las pendientes de validación comercial como por confirmar con proveedor. El punto de selección del ventilador se definió en condiciones de catálogo (3 840 m$^3$/h a 260 Pa, $\rho$ = 1.2 kg/m$^3$) aplicando el método del Ejemplo 3 de FE-1600 \cite{tcf_fe1600}, de modo que la selección exacta de tamaño y RPM se realice con el software del fabricante (CAPS de Greenheck o equivalente).
